\section{Discussion}

This section interprets the findings presented in the Results in relation
to the research objectives, existing literature, and the broader
theoretical and practical implications of the study.
 
% ------------------------------------------------------------
\subsection{Algorithm Selection for Mobile Route Optimization}
% ------------------------------------------------------------
 
The benchmarking results reveal a clear and consequential trade-off
between solution quality and runtime efficiency across the seven tested
algorithms. While GA achieved the highest accuracy with an average
relative error of only 0.10\%, its runtime scaled exponentially with
problem size, reaching 196.59 seconds on the \textit{berlin52} instance making
it impractical for deployment in a real-time mobile environment where
users expect immediate feedback upon requesting an optimized route.
This finding is consistent with the argument of \textcite{Wadi2025} that
heuristic models offering rapid execution with near-optimal accuracy are
substantially more effective for resource-constrained environments than
exact or computationally intensive methods that risk critical resource
exhaustion.
 
SA’s selection as the route optimization algorithm for \textit{Lakad} is
well-supported by its performance profile. With an average relative
error of 1.59\% and a maximum runtime of 1.90 seconds on the largest
tested instance. This accuracy margin is acceptable for tourism itinerary routing, where the difference between a near-optimal and optimal path represents minor distance variations unlikely to meaningfully affect user experience. EHO ranked 2nd under SAW due to its speed advantage, but its accuracy degraded
sharply on larger instances, reaching an average error of 11.27\% making it unsuitable for itineraries involving many stops. This highlights that raw speed alone is an insufficient criterion for algorithm selection in applied optimization contexts; the balance between both criteria, as operationalized through SAW, provides a more defensible selection methodology than either criterion in isolation.
 
The standardized TSPLIB benchmarks together with fixed algorithm parameters and 30
independent test runs per benchmark create a reliable testing method to choose optimization algorithms for
mobile resource-limited environments. The methodological contribution
of this research extends beyond \textit{Lakad} because the findings provide
real-world evidence which helps other applied systems that need to
perform real-time combinatorial optimization on mobile platforms.

% ------------------------------------------------------------
\subsection{Interpretation of System Acceptability (TAM)}
% ------------------------------------------------------------
 
The grand TAM mean score of 4.6719, interpreted as Highly Acceptable,
indicates that the target user population  found \textit{Lakad} to be both useful and easy to operate.
The high scores across Perceived Usefulness (4.7053), Perceived Ease of
Use (4.7067), and Attitude Towards Using (4.7200) suggest that the
system successfully addressed the core usability requirements identified
in the problem statement.

These results support the findings of \textcite{Yochum2020}, whose
development of AGAM demonstrated that multi-objective optimization
frameworks grounded in user preference modeling produce travel
recommendations that users find personally relevant and satisfying.
The high Perceived Usefulness score in particular reflects that users
recognized the value of receiving a personalized, constraint-aware
itinerary rather than a generic list of destinations, aligning with the documented demand for curated travel experiences reported by \textcite{SkiftState2022}, who found that 30 percent of travelers expressed increased preference for personalized planning services following the pandemic.
 
The Behavioral Intention to Use construct recorded the lowest mean score
(4.5556) and the highest standard deviation (0.4740) among all four
constructs. The item-level analysis provides a plausible explanation:
the Intention to Use item (BI1) received more 4-point than 5-point
responses, while the Future Use item (BI2) was predominantly rated 5.
This divergence suggests that while users are generally willing to
continue using the system, a subset expressed uncertainty about their
immediate intention, possibly because the system is currently limited
to Bulacan and single-day itineraries. This interpretation is further supported by the low standard
deviation in Perceived Usefulness (0.3435), indicating that uncertainty
about future use does not stem from doubts about the system’s functional value, but more likely from situational factors external to the application itself.
 
% ------------------------------------------------------------
\subsection{Interpretation of Software Quality (ISO/IEC 25010:2023)}
% ------------------------------------------------------------
 
The grand ISO/IEC 25010:2023 score of 4.7642 indicates that Lakad meets
professional software quality standards across all nine evaluated
constructs. The highest-scoring construct, Functional Suitability
(4.8636, $SD = 0.2447$), reflects near-unanimous expert consensus that
the system fulfills its core purpose. The lowest standard
deviation in this construct suggests that evaluators were most aligned
on the application’s capacity to perform its intended tasks, which
validates the design decision to scope the system narrowly around
itinerary generation rather than attempting to replicate the full
feature set of general-purpose travel platforms.
 
Security (4.8258) and Reliability (4.8182) ranked as the second and
third highest constructs, indicating that the combination of Supabase
for data management, encrypted storage of user data, and compliance
with the Data Privacy Act of 2012 (R.A. 10173) produced a system
that experts regarded as both trustworthy and stable. These scores
are particularly important for a system handling location data and
user-generated content, where data integrity and availability are
essential to user trust.
 
Compatibility received the lowest mean score (4.6591) and the highest
standard deviation (0.4973), which is directly attributable to the
system’s current Android-exclusive deployment. The elevated standard
deviation reflects genuine disagreement among evaluators about the
significance of this limitation, some experts likely weighted
cross-platform availability more heavily than others when assessing
the construct. The item-level analysis confirms this: the
Interoperability item accumulated the lowest frequency of 5-point
ratings within the Compatibility construct, while the Android-specific
performance items scored higher. Since \textit{Lakad} is developed in React Native, it open possibility for cross-platform deployment in the future with the existing codebase.
 
The lower scores on Analysability, Replaceability, and Modularity,
while still rated above average, warrant attention. The Replaceability
score suggests that evaluators recognized the competitive challenge
of positioning \textit{Lakad} against established platforms such as Google Maps
and TripAdvisor, which offer broader geographic coverage. This is
an inherent limitation of any localized application and should not
be interpreted as a functional deficiency, but rather as a market
positioning challenge that future development should address through
expanded geographic coverage and feature depth.
 
% ------------------------------------------------------------
\subsection{Practical and Policy Implications}
% ------------------------------------------------------------
 
For tourism practitioners and local government units, \textit{Lakad}
demonstrates that a low-cost, open-source algorithmic approach
can produce a functional and user-accepted destination promotion
tool for provinces that are underserved by commercial travel
platforms. The admin dashboard’s analytics provide PHACTO and local tourism operators with a data-driven basis for identifying which destinations receive
disproportionately low tourist traffic and adjusting promotional
strategies accordingly. This directly addresses the uneven
distribution of visitors across Bulacan’s municipalities
documented in the study’s problem context.

The system’s architecture with PHACTO-verified data,
open routing APIs, and a React Native mobile frontend, presents
a replicable template that other provincial tourism offices across
the Philippines could adopt with limited technical overhead.
The DOT’s National Tourism Development Plan 2023--2028, which
emphasizes the development of unique, localized travel experiences,
provides an explicit policy context within which systems like
\textit{Lakad} are not merely useful tools but strategic instruments
for sustainable provincial tourism development.
 
% ------------------------------------------------------------
\subsection{Limitations}
% ------------------------------------------------------------
 
Several limitations of this study should be acknowledged in
interpreting its findings. The \textit{Lakad} application has several scope and developmental limitations primarily driven by technical constraints, resource availability, and design choices. Geographically, the app is restricted to destinations within Bulacan and intentionally excludes restaurants from its Points of Interest (POIs) to avoid bias toward popular or larger establishments. The system caps itinerary generation at a maximum of 50 POIs and lacks public transit routing because of constraints within the integrated Mapbox API, which does not provide transit or multi-modal transportation data. Additionally, platform availability is limited exclusively to the Android Operating System, as developing for iOS required macOS hardware and an Apple Developer account that were not feasible for the study. Finally, the application’s functional scope is confined to generating single-day itineraries, lacking support for multi-day travel plans. And it does not integrate external booking services for accommodation reservations or ticket purchasing.

\subsection{Recommendations}
To further enhance \textit{Lakad}, future development should expand its database to include more destinations, accommodations, and pasalubong centers, while scaling its geographic reach to other provinces and developing an iOS version for wider accessibility. The system’s capabilities should be improved to support multi-day tour planning, integrate real-time traffic data with public transportation APIs, and immersive virtual exploration features. Future researchers should explore hybrid metaheuristic algorithms or different approach to potentially yield a much better routing and optimization. Lastly, migrating the system’s infrastructure to a self-hosted server rather than relying on external database and routing APIs.
